\documentclass[12pt]{article}
\usepackage[utf8]{inputenc}
\usepackage[spanish]{babel}
\usepackage{amsmath, amssymb, amsfonts}
\usepackage{graphicx}
\usepackage{geometry}
\usepackage{tikz}
\usepackage{xcolor}
\usepackage{booktabs}
\usepackage{tcolorbox}

\usetikzlibrary{shapes.geometric, arrows.meta, positioning, fit, backgrounds, calc}

\geometry{a4paper, margin=2.5cm}

\title{Formulación MDP del problema de portafolio con costos de transacción}
\author{}
\date{}

\begin{document}
\maketitle

Consideremos un conjunto finito de activos $I$, un horizonte discreto $t = 1, \dots, T$, y un capital inicial $x_1 > 0$. El problema de portafolio con costos de transacción proporcional puede modelarse como un Proceso de Decisión de Markov (MDP) en tiempo discreto.

\section{Definición del MDP}
Definimos el MDP como $\mathcal{M} = (\mathcal{S}, \mathcal{A}, \mathcal{P}, r, \gamma)$, donde:

\paragraph{Espacio de estados $\mathcal{S}$.} El estado al inicio del periodo $t$ se define como
$$s_t = (x_t, w_{t-1}, t),$$
donde:
\begin{itemize}
    \item $x_t \in \mathbb{R}^+$ es el capital disponible al inicio de $t$.
    \item $w_{t-1} = (w_{i,t-1})_{i \in I}$ son los pesos vigentes del portafolio al cierre de $t - 1$, con $w_{i,t-1} \geq 0$, $\sum_{i \in I} w_{i,t-1} = 1$.
    \item $t$ es el índice de tiempo (opcionalmente normalizado $t/T$ para efectos numéricos).
\end{itemize}
Por lo tanto, el espacio de estados es
$$\mathcal{S} = \left\{(x, w, t) : x \geq 0, w \in \Delta^{|I|-1}, t \in \{1, \dots, T\}\right\},$$
donde $\Delta^{|I|-1}$ denota el simplex de dimensión $|I| - 1$.

\paragraph{Espacio de acciones $\mathcal{A}$.} La acción en el periodo $t$ es la elección de un nuevo vector de pesos $w_t = (w_{i,t})_{i \in I}$:
$$a_t = w_t \in \mathcal{A},$$
sujeto a
$$w_{i,t} \geq 0,$$
$$\sum_{i \in I} w_{i,t} = 1.$$
En términos de RL continuo, $\mathcal{A}$ es también un simplex; en la implementación se trabaja con una parametrización continua que se proyecta al simplex.

\paragraph{Proceso de retornos.} Sea $R_t = (R_{i,t})_{i \in I}$ el vector de retornos en el intervalo $(t, t + 1]$. Supondremos que
$$R_t \sim D_t(\mu_t, \Sigma_t),$$
donde:
$$\mu_t = (\mu_{i,t})_{i \in I},$$
$$\Sigma_t = [\sigma_{ij,t}]_{i,j \in I}.$$
La distribución $D_t$ puede ser, por ejemplo, una distribución gaussiana multivariada o una mezcla condicionada a regímenes (bear/bull). Lo importante es que, dado $t$, la distribución de $R_t$ es conocida o estimada a partir de datos históricos.

\paragraph{Costos de transacción.} Dado el estado $(x_t, w_{t-1}, t)$ y la acción $w_t$, definimos el cambio de pesos
$$\Delta w_{i,t} = w_{i,t} - w_{i,t-1},$$
y la rotación total (en pesos del portafolio) como
$$\text{turn}_t = \sum_{i \in I} |\Delta w_{i,t}|.$$
Sea $c_i \geq 0$ la tasa de costo de transacción por unidad de rotación del activo $i$. El costo monetario en el periodo $t$ es
$$C_t(x_t, w_{t-1}, w_t) = x_t \sum_{i \in I} c_i |\Delta w_{i,t}|.$$

\paragraph{Dinámica de transición $\mathcal{P}$.} Dado el estado actual $s_t = (x_t, w_{t-1}, t)$, la acción $a_t = w_t$ y el retorno aleatorio $R_t$, definimos:
\begin{enumerate}
    \item Retorno del portafolio en el periodo $t$:
    $$r^{\text{port}}_t = \sum_{i \in I} w_{i,t} R_{i,t} = w_t^\top R_t.$$
    \item Capital al inicio de $t + 1$:
    $$x_{t+1} = x_t \left(1 + w_t^\top R_t - \sum_{i \in I} c_i |\Delta w_{i,t}|\right).$$
    \item Pesos para el siguiente estado: $w_t$ pasa a ser la componente de pesos del siguiente estado.
    \item Índice temporal: $t \to t + 1$.
\end{enumerate}
Por tanto, el kernel de transición queda implícitamente definido por:
$$S_{t+1} = (x_{t+1}, w_t, t + 1) = F(S_t, A_t, R_t),$$
con $R_t \sim D_t(\mu_t, \Sigma_t)$.

\paragraph{Función de recompensa $r(s, a)$.} Existen dos enfoques naturales y compatibles con el modelo de media--varianza con costos:
\begin{enumerate}
    \item Reward tipo media--varianza (determinístico):
    $$r_t(s_t, a_t) = \sum_{i \in I} w_{i,t} \mu_{i,t} - \lambda \sum_{i,j \in I} w_{i,t} w_{j,t} \sigma_{ij,t} - \sum_{i \in I} c_i |\Delta w_{i,t}|,$$
    donde $\lambda > 0$ es el parámetro de aversión al riesgo. Con $\gamma = 1$ y horizonte finito, el funcional $\sum_{t=1}^T r_t(s_t, a_t)$ reproduce la función objetivo de media--varianza con costos.
    \item Reward monetario (realizado):
    $$\tilde{r}_t(s_t, a_t, R_t) = x_{t+1} - x_t = x_t \left( w_t^\top R_t - \sum_{i \in I} c_i |\Delta w_{i,t}| \right),$$
    y el agente observa realizaciones de $\tilde{r}_t$ a partir de los retornos históricos o simulados.
\end{enumerate}
En ambos casos, el objetivo del agente es maximizar la suma descontada de recompensas:
$$J(\pi) = \mathbb{E}_\pi \left[ \sum_{t=1}^T \gamma^{t-1} r_t(S_t, A_t) \right],$$
para una política $\pi$.

\subsection*{1.2 Políticas y funciones de valor}
Una política estocástica se define como
$$\pi(w | x, w_{t-1}, t) = \Pr\{A_t = w | X_t = x, W_{t-1} = w_{t-1}, t\},$$
mientras que una política determinista se modela como
$$\mu : \mathcal{S} \to \mathcal{A}, \quad A_t = \mu(S_t).$$
Dada una política $\pi$, la función valor de estado es:
$$V^\pi(x, w_{t-1}, t) = \mathbb{E}_\pi \left[ \sum_{\tau=t}^T \gamma^{\tau-t} r_\tau(S_\tau, A_\tau) \Big| S_t = (x, w_{t-1}, t) \right].$$
La ecuación de Bellman para $V^\pi$ (con horizonte finito y $t < T$) es:
$$V^\pi(s_t) = \mathbb{E}_{A_t \sim \pi} \left[ r_t(s_t, A_t) + \gamma \mathbb{E}_{R_t} [V^\pi(S_{t+1})] \right].$$
Análogamente, la función valor acción--estado $Q^\pi$ es:
$$Q^\pi(s_t, a_t) = \mathbb{E} \left[ \sum_{\tau=t}^T \gamma^{\tau-t} r_\tau(S_\tau, A_\tau) \Big| S_t = s_t, A_t = a_t \right],$$
y satisface la ecuación de Bellman correspondiente.

Adicionalmente, para evaluar la calidad relativa de una acción respecto al valor promedio del estado, definimos la \textbf{Función de Ventaja (Advantage)}:
$$A^\pi(s_t, a_t) = Q^\pi(s_t, a_t) - V^\pi(s_t).$$
Empíricamente, en la arquitectura de algoritmos tipo Actor-Crítico (como PPO), esta ventaja se estima mediante la diferencia temporal:
$$A_t \approx r_t + \gamma V^\pi(S_{t+1}) - V^\pi(S_t).$$

Las funciones óptimas se definen por
$$V^*(s) = \max_\pi V^\pi(s),$$
$$Q^*(s, a) = \max_\pi Q^\pi(s, a),$$
y una política óptima $\pi^*$ cumple
$$\pi^*(s) \in \arg\max_a Q^*(s, a).$$

\section{Metodología de implementación DRL (PPO y SAC)}
A continuación se presenta una metodología paso a paso para implementar este MDP de portafolio con costos de transacción usando aprendizaje por refuerzo profundo (DRL) con stable baselines3, PyTorch y los algoritmos PPO y SAC. La metodología se organiza en etapas coherentes desde la construcción del entorno hasta la comparación con la solución de referencia obtenida vía GAMS en un esquema de rolling horizon.

\subsection{Etapa 1: Preparación de datos y estimación de parámetros}
\begin{enumerate}
    \item \textbf{Datos de retornos.} Reunir series históricas semanales (o en la frecuencia de trabajo) de retornos $R_{i,t}$ para todos los activos $i \in I$ en el horizonte de entrenamiento.
    \item \textbf{Estimación de momentos por periodo.} Para cada periodo $t$, estimar:
    $$\mu_t = \mathbb{E}[R_t],$$
    $$\Sigma_t = \text{Cov}(R_t),$$
    utilizando:
    \begin{itemize}
        \item Un modelo condicional con regímenes (bear/bull) y mezcla de probabilidades, o
        \item Ventanas móviles de datos históricos.
    \end{itemize}
    Estos objetos $\mu_t$ y $\Sigma_t$ son equivalentes a los parámetros $\mu_{mix}$ y $\Sigma_{mix}$ utilizados en el modelo GAMS.
    \item \textbf{Partición train/validation/test.} Dividir la serie de tiempo total en tres segmentos ordenados:
    \begin{itemize}
        \item Train: periodos iniciales para aprendizaje del agente.
        \item Validation: tramo intermedio para selección de hiperparámetros y control de sobreajuste.
        \item Test: tramo final para evaluar error de generalización y desempeño out-of-sample.
    \end{itemize}
\end{enumerate}

\subsection{Etapa 2: Definición del entorno MDP (Gym-like)}
Implementar un entorno en la interfaz tipo Gym, por ejemplo \texttt{PortfolioEnv}, con los métodos \texttt{reset()} y \texttt{step(action)}:
\begin{enumerate}
    \item \textbf{Observaciones.} Definir el vector de observación en $t$ como
    $$o_t = \left[ x_t, w_{t-1}, t/T, \text{features}_t \right],$$
    donde $\text{features}_t$ se define explícitamente como los parámetros normalizados del modelo GAMS:
    $$\text{features}_t = \left[ \text{norm\_}\mu_{\text{mix}}(i,t), \text{norm\_var\_mix}(i,t), \text{norm\_cov\_mix}(t) \right].$$
    Normalizar o escalar las observaciones para estabilidad numérica.
    \item \textbf{Acciones.} El agente devuelve un vector continuo $\tilde{a}_t \in \mathbb{R}^{|I|}$ que se transforma a pesos factibles:
    $$w_{i,t} = \frac{\max(0, \tilde{a}_{i,t})}{\sum_{j \in I} \max(0, \tilde{a}_{j,t})},$$
    garantizando $w_{i,t} \geq 0$ y $\sum_i w_{i,t} = 1$.
    \item \textbf{Step.} Dado el estado interno $(x_t, w_{t-1}, t)$ y la acción $w_t$:
    \begin{enumerate}
        \item[(a)] Generar o leer $R_t$:
        \begin{itemize}
            \item Opción simulación: muestrear $R_t \sim D_t(\mu_t, \Sigma_t)$.
            \item Opción histórica: tomar el retorno observado en $t$.
        \end{itemize}
        \item[(b)] Calcular el costo de transacción:
        $$C_t = x_t \sum_{i \in I} c_i |w_{i,t} - w_{i,t-1}|.$$
        \item[(c)] Actualizar el capital:
        $$x_{t+1} = x_t \left(1 + w_t^\top R_t - \sum_i c_i |w_{i,t} - w_{i,t-1}|\right).$$
        \item[(d)] Actualizar pesos: $w^{\text{next}}_t = w_t$.
        \item[(e)] Definir el reward $r_t$ según la opción elegida:
        $$r_t = \sum_i w_{i,t} \mu_{i,t} - \lambda \sum_{i,j} w_{i,t} w_{j,t} \sigma_{ij,t} - \sum_i c_i |\Delta w_{i,t}| \quad \text{(media--varianza)}$$
        o $x_{t+1} - x_t$ (monetario realizado).
        \item[(f)] Construir la observación $o_{t+1}$ y marcar \texttt{done = True} cuando $t = T$.
    \end{enumerate}
    \item \textbf{Reset.} En \texttt{reset()}, inicializar:
    $$x_1 = x_{\text{init}}, \quad w_0 = w_{\text{init}}, \quad t = 1,$$
    y devolver la observación inicial $o_1$.
\end{enumerate}

\vspace{1cm}
\begin{center}
\textbf{\Large Diagrama de Arquitectura PPO}
\end{center}
\vspace{0.5cm}

% -------------------------
% DIAGRAMA PPO TIKZ AQUI
% -------------------------
\begin{center}
\resizebox{\textwidth}{!}{
\begin{tikzpicture}[
    >=Latex,
    node distance=1.5cm and 2cm,
    font=\sffamily,
    % Estilos de cajas
    agentbox/.style={rectangle, rounded corners, draw=blue!60, fill=blue!5, very thick, align=center, inner sep=10pt},
    envbox/.style={rectangle, rounded corners, draw=green!60!black, fill=green!5, very thick, align=center, inner sep=10pt},
    mathbox/.style={rectangle, draw=gray!50, fill=white, align=center, inner sep=6pt},
    arrowlabel/.style={fill=white, text=black, font=\footnotesize, align=center}
]

% AGENTE PPO (Izquierda)
\node[agentbox] (actor) {
    \textbf{Red Actor (Política $\pi_\theta$)}\\[1ex]
    Input: $O_t$\\[1ex]
    Salida: $a_t$ (Acción Continua)
};



\node[agentbox, below=1.5cm of muestreo] (critic) {
    \textbf{Red Crítico (Valor $V_\phi$)}\\[1ex]
    Input: $O_t$\\[1ex]
    Salida: $V(O_t)$
};

\node[mathbox, below=0.8cm of critic] (ventaja) {
    \textbf{Cálculo de Ventaja ($A_t$)}
};

\node[mathbox, below=0.8cm of ventaja] (loss) {
    \textbf{Optimizador PPO}
};

% Contenedor del Agente
\begin{scope}[on background layer]
    \node[rectangle, rounded corners, draw=blue!80, fill=blue!2, dashed, very thick, fit=(actor) (muestreo) (critic) (ventaja) (loss), inner sep=15pt, label={[font=\large\bfseries, text=blue!80]above:Cerebro PPO (Agente On-Policy)}] (agent_bg) {};
\end{scope}

% Conexiones internas Agente

\draw[->, thick, blue] (critic) -- (ventaja) node[midway, right, font=\footnotesize] {Predicción};
\draw[->, thick, blue, dashed] (ventaja) -- (loss) node[midway, right, font=\footnotesize] {Optimiza pesos};

% ENTORNO (Derecha)
\node[envbox, right=5cm of muestreo] (proyeccion) {
    \textbf{Proyección Símplex ($w_t$)}
};

\node[envbox, below=0.8cm of proyeccion] (dinamica) {
    \textbf{Evolución Capital ($x_{t+1}$)}
};

\node[envbox, below=0.8cm of dinamica] (recompensa) {
    \textbf{Reward Media-Varianza ($r_t$)}
};

% Contenedor del Entorno
\begin{scope}[on background layer]
    \node[rectangle, rounded corners, draw=green!60!black, fill=green!2, dashed, very thick, fit=(proyeccion) (dinamica) (recompensa), inner sep=15pt, label={[font=\large\bfseries, text=green!60!black]above:PortfolioEnv (Mercado)}] (env_bg) {};
\end{scope}

% Conexiones internas Entorno
\draw[->, thick, green!60!black] (proyeccion) -- (dinamica);
\draw[->, thick, green!60!black] (dinamica) -- (recompensa);

% INTERACCIONES (Ciclo RL)
\draw[->, line width=1.5pt, orange] (actor) -- (proyeccion) node[midway, arrowlabel] {Acción Continua\\$a_t \in \mathbb{R}^{|I|}$};

\draw[->, line width=1.5pt, red] (recompensa.west) -- (ventaja.east) node[midway, fill=white, font=\footnotesize, align=center] {Señal $r_t$};

\coordinate (obs_out) at ($(recompensa.south) + (0, -0.5)$);
\draw[line width=1.5pt, purple] (recompensa.south) -- (obs_out);
\draw[line width=1.5pt, purple] (obs_out) -| ($(agent_bg.west) + (-1.5, 0)$) |- (actor.west) 
    node[pos=0.10, arrowlabel, xshift=2.5cm] {\textbf{Vector de Observación} $O_{t+1}$\\$O_{t+1} = \left[x_{t+1}, w_t, \frac{t}{T}, \text{norm\_}\mu_{\text{mix}}, \text{norm\_var}_{\text{mix}}, \text{norm\_cov}_{\text{mix}}\right]$};

\draw[->, line width=1.5pt, purple] ($(agent_bg.west) + (-1.5, 0)$) |- (critic.west);

\end{tikzpicture}
}
\end{center}

\subsection{Etapa 3: Definición de agentes y arquitecturas (PPO y SAC)}
\textbf{Espacio de observación y acción.} Configurar en stable baselines3:
\begin{itemize}
    \item Box continuo para el espacio de observación, con límites apropiados.
    \item Box continuo para el espacio de acción, con límites, por ejemplo, en $[-1, 1]$ y posterior proyección al simplex como se indicó.
\end{itemize}

\textbf{PPO (on-policy, policy gradient).}
\begin{itemize}
    \item Usar PPO con política tipo MlpPolicy.
    \item Arquitectura sugerida: 2–3 capas densas intermedias con activaciones no lineales (ReLU o Tanh), tamaño moderado (p.ej., 64–128 neuronas) para evitar sobreparametrizar.
    \item Hiperparámetros a calibrar: tamaño de batch, longitud de rollouts, coeficiente de entropía, clip range, tasa de aprendizaje, coeficiente de valor.
\end{itemize}

\textbf{SAC (off-policy, actor–critic maximum-entropy).}
\begin{itemize}
    \item Usar SAC con política MlpPolicy y acción continua.
    \item Arquitecturas similares para actor y críticos (2–3 capas densas).
    \item Hiperparámetros clave: tasa de aprendizaje, tau para target networks, temperatura de entropía (fija o automática).
\end{itemize}


\textbf{Funciones de Pérdida (Loss) para Backpropagation.}
Para materializar el aprendizaje, los algoritmos optimizan las siguientes funciones de pérdida mediante gradiente descendente:

\textbf{Loss de PPO:}
El agente PPO actualiza su política $\pi_\theta$ maximizando una función objetivo recortada (Clipped Surrogate Objective) para evitar actualizaciones destructivas:
$$L^{CLIP}(\theta) = \hat{\mathbb{E}}_t \left[ \min\left( \rho_t(\theta) \hat{A}_t, \text{clip}(\rho_t(\theta), 1-\epsilon, 1+\epsilon) \hat{A}_t \right) \right]$$
donde $\rho_t(\theta) = \frac{\pi_\theta(a_t | s_t)}{\pi_{\theta_{old}}(a_t | s_t)}$ es la razón de probabilidad y $\epsilon$ es el rango de recorte. Simultáneamente, el Crítico $V_\phi$ se actualiza minimizando el error cuadrático medio respecto a los retornos calculados:
$$L^{VF}(\phi) = \hat{\mathbb{E}}_t \left[ \left( V_\phi(s_t) - V_t^{\text{target}} \right)^2 \right]$$

\textbf{Loss de SAC:}
En SAC, los parámetros de las redes críticas gemelas $Q_{\phi_i}$ se optimizan minimizando el Error Cuadrático Medio de la ecuación de Bellman \textit{Soft}:
$$L_Q(\phi_i) = \mathbb{E}_{(s_t, a_t) \sim \mathcal{D}} \left[ \left( Q_{\phi_i}(s_t, a_t) - \left( r_t + \gamma \left( \min_{j=1,2} Q_{\phi_{\text{target}, j}}(s_{t+1}, \tilde{a}_{t+1}) - \alpha \log \pi_\theta(\tilde{a}_{t+1} | s_{t+1}) \right) \right) \right)^2 \right]$$
Por su parte, el Actor $\pi_\theta$ actualiza sus pesos minimizando la pérdida de política basada en la divergencia Kullback-Leibler, lo que se traduce en:
$$L_\pi(\theta) = \mathbb{E}_{s_t \sim \mathcal{D}, \tilde{a}_t \sim \pi_\theta} \left[ \alpha \log \pi_\theta(\tilde{a}_t | s_t) - \min_{j=1,2} Q_{\phi_j}(s_t, \tilde{a}_t) \right]$$

\textbf{Capacidad de los agentes y error de generalización.}
El objetivo es minimizar el error de generalización (diferencia entre desempeño en entrenamiento y en test). Para ello:
\begin{enumerate}
    \item \textbf{Control de capacidad.} Ajustar la profundidad y el ancho de las redes para que capturen patrones relevantes sin sobreajuste evidente.
    \item \textbf{Regularización implícita.} 
    \begin{itemize}
        \item En PPO: controlar clip range y tasa de aprendizaje para evitar updates agresivos.
        \item En SAC: ajustar la temperatura de entropía para evitar políticas excesivamente deterministas durante el entrenamiento.
    \end{itemize}
    \item \textbf{Evaluación en validación.} Durante el entrenamiento, usar un evaluation callback que mida periódicamente:
    $$\hat{J}_{\text{val}}(\theta) = \frac{1}{N_{ep}} \sum_{e=1}^{N_{ep}} (\text{retorno acumulado en episodios de validación}),$$
    y seleccionar el modelo con mejor desempeño en validación.
    \item \textbf{Walk-forward y múltiples semillas.} Repetir el entrenamiento con distintas semillas iniciales y, opcionalmente, distintos sub-períodos de entrenamiento para medir robustez y estabilidad.
\end{enumerate}


\vspace{0.5cm}
\textbf{Definiciones para Arquitectura SAC.} 
Para el agente \textbf{SAC (Soft Actor-Critic)}, el objetivo se expande para maximizar tanto el retorno esperado como la entropía de la política. El proceso incluye componentes específicos para su naturaleza off-policy:
\begin{itemize}
    \item \textbf{Replay Buffer ($\mathcal{D}$):} Memoria donde se almacenan las transiciones observadas $(O_t, a_t, r_t, O_{t+1})$. Permite reutilizar experiencia pasada.
    \item \textbf{Doble Crítico Soft ($Q_{\phi_1}, Q_{\phi_2}$):} Dos redes neuronales que estiman el valor acción-estado continuo $Q(O_t, a_t)$. Se utiliza el mínimo de ambas estimaciones para evitar el sesgo de sobreestimación del valor.
    \item \textbf{Optimizador SAC ($\alpha$, $\tau$):} Incluye el cálculo y optimización de las \textbf{Redes Target} mediante una media móvil exponencial con tasa $\tau$, y la gestión del \textbf{Coeficiente de Entropía ($\alpha$)}, el cual controla la temperatura de exploración.
\end{itemize}

A continuación se presenta la arquitectura conceptual para el agente SAC, la cual interactúa con el mismo entorno definido previamente:

\vspace{0.5cm}
\begin{center}
\textbf{\Large Diagrama de Arquitectura SAC}
\end{center}
\vspace{0.3cm}

\begin{center}
\resizebox{\textwidth}{!}{
\begin{tikzpicture}[
    >=Latex,
    node distance=1.5cm and 2cm,
    font=\sffamily,
    % Estilos de cajas
    agentbox/.style={rectangle, rounded corners, draw=blue!60, fill=blue!5, very thick, align=center, inner sep=10pt},
    envbox/.style={rectangle, rounded corners, draw=green!60!black, fill=green!5, very thick, align=center, inner sep=10pt},
    mathbox/.style={rectangle, draw=gray!50, fill=white, align=center, inner sep=6pt},
    arrowlabel/.style={fill=white, text=black, font=\footnotesize, align=center}
]

% AGENTE SAC (Izquierda)
\node[agentbox] (actor) {
    \textbf{Red Actor (Política Determinista $\mu_\theta$)}\\[1ex]
    Input: $O_t$\\[1ex]
    Salida: $a_t$ (Acción Continua)
};



\node[mathbox, below=1cm of actor] (buffer) {
    \textbf{Replay Buffer ($\mathcal{D}$)}
};

\node[agentbox, below=1cm of buffer] (critic) {
    \textbf{Doble Crítico Soft ($Q_{\phi_1}, Q_{\phi_2}$)}\\[1ex]
    Input: $O_t, a_t$\\[1ex]
    Salida: $Q(O_t, a_t)$
};

\node[mathbox, below=0.8cm of critic] (loss) {
    \textbf{Optimizador SAC ($\alpha$, $\tau$)}
};

% Contenedor del Agente
\begin{scope}[on background layer]
    \node[rectangle, rounded corners, draw=blue!80, fill=blue!2, dashed, very thick, fit=(actor) (buffer) (critic) (loss), inner sep=15pt, label={[font=\large\bfseries, text=blue!80]above:Cerebro SAC (Agente Off-Policy)}] (agent_bg) {};
\end{scope}

% Conexiones internas Agente

\draw[->, thick, blue, dashed] (buffer) -- (critic) node[midway, right, font=\footnotesize] {Minibatch};
\draw[->, thick, blue, dashed] (critic) -- (loss) node[midway, right, font=\footnotesize] {Optimiza pesos};

% ENTORNO (Derecha) - Alineado con SAC
\node[envbox, right=6cm of actor, yshift=-1.5cm] (proyeccion) {
    \textbf{Proyección Símplex ($w_t$)}
};

\node[envbox, below=1.2cm of proyeccion] (dinamica) {
    \textbf{Evolución Capital ($x_{t+1}$)}
};

\node[envbox, below=1.2cm of dinamica] (recompensa) {
    \textbf{Reward Media-Varianza ($r_t$)}
};

% Contenedor del Entorno
\begin{scope}[on background layer]
    \node[rectangle, rounded corners, draw=green!60!black, fill=green!2, dashed, very thick, fit=(proyeccion) (dinamica) (recompensa), inner sep=15pt, label={[font=\large\bfseries, text=green!60!black]above:PortfolioEnv (Mercado)}] (env_bg) {};
\end{scope}

% Conexiones internas Entorno
\draw[->, thick, green!60!black] (proyeccion) -- (dinamica);
\draw[->, thick, green!60!black] (dinamica) -- (recompensa);

% INTERACCIONES (Ciclo RL)
\draw[->, line width=1.5pt, orange] (actor) -- (proyeccion) node[midway, arrowlabel] {Acción Continua\\$a_t \in \mathbb{R}^{|I|}$};

% La señal de recompensa va al buffer
\draw[->, line width=1.5pt, red] (recompensa.west) -- (buffer.east) node[midway, fill=white, font=\footnotesize, align=center] {Transición\\$(O_t, a_t, r_t, O_{t+1})$};

\coordinate (obs_out) at ($(recompensa.south) + (0, -0.5)$);
\draw[line width=1.5pt, purple] (recompensa.south) -- (obs_out);
\draw[line width=1.5pt, purple] (obs_out) -| ($(agent_bg.west) + (-1.5, 0)$) |- (actor.west) 
    node[pos=0.10, arrowlabel, xshift=3.5cm] {\textbf{Vector de Observación} $O_{t+1}$\\$O_{t+1} = \left[x_{t+1}, w_t, \frac{t}{T}, \text{norm\_}\mu_{\text{mix}}, \text{norm\_var}_{\text{mix}}, \text{norm\_cov}_{\text{mix}}\right]$};

\draw[->, line width=1.5pt, purple] ($(agent_bg.west) + (-1.5, 0)$) |- (critic.west);

\end{tikzpicture}
}
\end{center}
\vspace{1cm}

\subsection{Etapa 4: Diseño de pruebas de comparación}
Se propone un esquema de comparación con las siguientes referencias:
\begin{enumerate}
    \item \textbf{Base determinista (GAMS, media–varianza con costos).}
    \begin{itemize}
        \item Resolver el modelo media–varianza con costos de transacción en GAMS para ventanas móviles de longitud $H$ (por ejemplo, 26–52 semanas) usando las estimaciones $\mu_t$ y $\Sigma_t$ disponibles hasta el tiempo actual.
        \item En cada fecha $t$:
        \begin{enumerate}
            \item[(a)] Construir una ventana rolling $[t, t+H-1]$ con parámetros pronosticados.
            \item[(b)] Resolver el problema determinista en esa ventana (maximización de media–varianza con costos).
            \item[(c)] Aplicar sólo la primera decisión $w^{\text{GAMS}}_t$.
            \item[(d)] Avanzar al siguiente periodo $t + 1$ y repetir (esquema de rolling horizon).
        \end{enumerate}
        \item Evolucionar el capital $x^{\text{GAMS}}_t$ con los retornos realizados y registrar la trayectoria.
    \end{itemize}
    \item \textbf{Agentes DRL (PPO y SAC).}
    \begin{itemize}
        \item Para cada agente entrenado (PPO y SAC), fijar los parámetros finales $\theta^\star$ y evaluar en el periodo de test:
        $$x^{\text{PPO}}_t, \quad x^{\text{SAC}}_t, \quad t \in \text{test},$$
        usando retornos realizados y el mismo esquema de costos.
    \end{itemize}
    \item \textbf{Benchmarks adicionales.}
    \begin{itemize}
        \item Portafolio buy-and-hold (por ejemplo, 50/50).
        \item Portafolio 50/50 con rebalanceo periódico fijo.
    \end{itemize}
\end{enumerate}

\subsection{Etapa 5: Métricas de evaluación y comparación}
Para comparar GAMS vs. PPO vs. SAC vs. benchmarks se propone:
\begin{itemize}
    \item Retorno acumulado:
    $$\text{RetAcum} = \frac{x_T}{x_1} - 1.$$
    \item Volatilidad del portafolio: desviación estándar de los retornos de portafolio $r^{\text{port}}_t$.
    \item Medidas riesgo–retorno: ratio de Sharpe, semivarianza, drawdowns máximos.
    \item Rotación y costos: promedio de $\text{turn}_t$ y costo acumulado de transacción.
    \item Error de generalización: diferencia entre métricas en train/validation y test para cada agente, y comparación del agente DRL contra la solución rolling de GAMS:
    $$\text{Regret}_{\text{GAMS}}^{\text{DRL}} = \text{RetAcum}_{\text{GAMS}} - \text{RetAcum}_{\text{DRL}},$$
    evaluado en el periodo de test.
\end{itemize}

\subsection{Etapa 6: Síntesis}
La metodología propuesta permite:
\begin{enumerate}
    \item Formular el problema de portafolio con costos de transacción como un MDP bien definido.
    \item Entrenar agentes DRL (PPO y SAC) que aprenden políticas $\pi_\theta$ o $\mu_\theta$ a partir de datos históricos o simulados.
    \item Controlar la capacidad de los agentes y el error de generalización mediante particiones temporales, callbacks de evaluación y experimentos con múltiples semillas.
    \item Comparar de forma rigurosa el desempeño de los agentes DRL contra la solución determinista GAMS en un esquema de rolling horizon, y contra benchmarks simples, tanto en términos de retorno esperado como de riesgo y costos de transacción.
\end{enumerate}



\clearpage
\section{Apéndice: Preparación de Datos y Estimación de Parámetros}

Este documento detalla la ejecución técnica de la \textbf{Etapa 1} definida en la Memoria de Referencia (Documento 02), adaptándola a los activos SPX y CMC200 para un horizonte de $T=163$ semanas.

\section{Datos Fuente y Horizonte Temporal}

El proyecto utiliza series históricas semanales para $T=163$ periodos. Los insumos base son:
\begin{itemize}
    \item \textbf{Retornos Netos ($R_{i,t}$):} Observaciones semanales de SPX y CMC200.
    \item \textbf{Probabilidades HMM ($P_{i,k,t}$):} Señales exógenas de régimen (Bear/Bull).
\end{itemize}

\section{Partición del Dataset (Rigor DRL)}

Siguiendo la metodología \textit{walk-forward} de la Biblia, dividimos las 163 semanas para garantizar la capacidad de generalización:

\begin{table}[h!]
\centering
\begin{tabular}{@{}lll@{}}
\toprule
\textbf{Conjunto} & \textbf{Rango (semanas)} & \textbf{Propósito} \\ \midrule
Entrenamiento (Train) & 1 - 114 (70\%) & Aprendizaje de la política. \\
Validación (Val) & 115 - 138 (15\%) & Ajuste de hiperparámetros. \\
Prueba (Test) & 139 - 163 (15\%) & Evaluación out-of-sample vs GAMS. \\ \bottomrule
\end{tabular}
\end{table}

\section{Estimación de Parámetros Mix}

Para cada periodo $t$, se calculan las expectativas de mercado mediante una ventana expansiva ($1 \dots t$).

\subsection{Estimación de Momentos por Régimen}
Para cada activo $i$ bajo el régimen $k$, los momentos se estiman ponderando cada observación histórica por su probabilidad de pertenencia al régimen:

\begin{equation}
    \hat{\mu}_{i,k,t} = \frac{\sum_{m=1}^t P_{i,k,m} R_{i,m}}{\sum_{m=1}^t P_{i,k,m}}
\end{equation}

\begin{equation}
    \hat{\sigma}_{i,j,k,t} = \frac{\sum_{m=1}^t P_{i,k,m} (R_{i,m} - \hat{\mu}_{i,k,t})(R_{j,m} - \hat{\mu}_{j,k,t})}{\sum_{m=1}^t P_{i,k,m}}
\end{equation}

\subsection{Construcción de Variables Mix}
Integramos la incertidumbre del régimen en una única señal de media y riesgo:
\begin{equation}
    \mu_{mix, i,t} = \sum_{k} P_{i,k,t} \cdot \hat{\mu}_{i,k,t}
\end{equation}
\begin{equation}
    \sigma_{mix, i,j,t} = \sum_{k} P_{i,k,t} \cdot \hat{\sigma}_{i,j,k,t}
\end{equation}

Estas variables simplifican la tarea del agente al entregarle la información ya integrada.

\section{Normalización Z-Score de Observaciones}

Aplicamos estandarización expansiva para asegurar estabilidad numérica. Para cada variable $x_{mix, t}$, calculamos su versión normalizada:
\begin{equation}
    z(x)_{mix, t} = \frac{x_{mix, t} - \mu(x)_{mix, t}}{\sigma(x)_{mix, t}}
\end{equation}

\section{Parámetros del Modelo y Configuración Base}

Para asegurar la comparabilidad con el modelo GAMS original, se registran las siguientes constantes:

\begin{itemize}
    \item \textbf{Capital Inicial ($x_0$):} 10,000 USD.
    \item \textbf{Pesos Iniciales ($w_0$):} 50\% SPX / 50\% CMC200.
    \item \textbf{Costos de Transacción ($c_i$):} SPX ($0.5\%$), CMC200 ($1.0\%$).
    \item \textbf{Aversión al Riesgo ($\lambda$):}
    \begin{itemize}
        \item Grilla GAMS: $\{0.05, 0.10, 0.25, 0.50, 1.00\}$.
        \item \textbf{Selección DRL:} $\lambda = 0.10$.
    \end{itemize}
\end{itemize}

\end{document}